\chapter{Calculations for explanatory preferences}
\label{annex:p19-explanatory-preferences}

\begin{annexitem}{annex:p19.explanation-logical-menu}{thm:p19-explanation-logical-menu}{The candidate menu as finite abduction}
\begin{proof}
Every subset \(A\subseteq\{D_1,D_2,D_3\}\) is consistent with
the two biconditionals: set all three disease atoms true and
both symptoms true. If \(D_1\in A\), both symptoms follow.
If \(D_2,D_3\in A\), the first supplies \(S_1\) and the
second supplies \(S_2\).

If neither condition holds, set \(D_1\) false and set each
of \(D_2,D_3\) true exactly when it belongs to \(A\).
At least one of these last two atoms is false.
Define each symptom by its corresponding disjunction.
This valuation satisfies \(T\cup A\) but makes at least
one observation false. Such an \(A\) is not an explanation.
Consequently the explanations are precisely the subsets
containing \(D_1\) or containing both \(D_2,D_3\).
Their only inclusion-minimal members are \(\{D_1\}\) and
\(\{D_2,D_3\}\); the former alone has least cardinality.
This uses exactly the finite criterion of
\cref{thm:p16-abduction-semantics}.

For the extra prevention rule, \(D_2,D_3\) would imply
\(\neg S_2,S_2\), respectively, so consistency fails.
Without that added rule, the valuation
\(D_1=\mathbf F,D_2=D_3=S_1=S_2=\mathbf T\) satisfies the
intended theory. This is a complete countervaluation to
identifying absence of a causal link with prevention by
the corresponding cause.
\end{proof}
\end{annexitem}

\begin{annexitem}{annex:p19.explanation-probability}{thm:p19-explanation-probability}{Inclusive events and dependence}
\begin{proof}
Inclusion--exclusion gives
\(\mathbb P(E)=p_1+j-t\). Both explanatory events are
subsets of \(E\). Dividing their probabilities by its
positive probability proves both conditional formulas.
Taking their ratio when \(j>0\) gives \(R=p_1/j\).
If \(p_1>0\), equality of
\(p_1/(p_1+j-t)\) and \(p_1/(p_1+j)\) is equivalent to
equality of the positive denominators, hence to \(t=0\).
Their posterior sum is
\((p_1+j)/(p_1+j-t)=1+t/\mathbb P(E)\).
Thus overlapping posteriors need not sum to one.

Independence of \(D_2,D_3\) gives
\(j=\mathbb P(D_2)\mathbb P(D_3)\) by its definition.
Marginal probabilities do not supply that premise.
Substituting \(p_1=50/750\), \(\mathbb P(D_2)=n_2/750\)
and \(\mathbb P(D_3)=n_3/750\), with \(n_2,n_3>0\),
gives \(R=37500/(n_2n_3)\).
The displayed \(250,220\) row therefore gives
\(37500/55000=15/22\), proving the table calculation.
For the independent fictional model take
\(\Omega=\{0,1\}^3\), with atom mass
\[
 (.1)^{d_1}(.9)^{1-d_1}
 (.4)^{d_2}(.6)^{1-d_2}
 (.4)^{d_3}(.6)^{1-d_3}.
\]
The masses are nonnegative and sum to one by multiplying
the three marginal sums. Define \(D_i=\{d_i=1\}\).
Then \(j=.16\), \(t=.016\), and
\(\mathbb P(E)=.244=61/250\).
The numerator masses are \(25/250\) and \(40/250\),
giving \(25/61\) and \(40/61\).
Their ratio is \(5/8\), and normalization over the two
numerator masses gives \(25/(25+40)=5/13\).

For two different joint laws with the same marginals,
give \((D_2,D_3)\) these masses:
\[
\begin{array}{c|rrrr}
(d_2,d_3)&(0,0)&(1,0)&(0,1)&(1,1)\\\hline
\text{Exclusive}&.2&.4&.4&0\\
\text{Identical}&.6&0&0&.4
\end{array}
\]
Each row sums to one and gives marginal probability .4
for each disease. Independently attach a Bernoulli(.1)
coordinate for \(D_1\), thereby giving a complete
probability law on the same eight outcomes.
For the first row \(j=t=0\), giving posteriors 1 and 0.
For the second, \(j=.4,t=.04\), so
\(\mathbb P(E)=.46=23/50\), giving \(5/23,20/23\).
All the claimed dependence comparisons are therefore
realized by fully specified finite laws.

For the equal-prevalence statement, independent \(D_2,D_3\)
give inclusive probabilities \(p\) and \(p^2\).
Their difference \(p(1-p)\) is positive for \(0<p<1\).
Under mutual independence, the exclusive configurations
\(D_1\cap D_2^c\cap D_3^c\) and
\(D_1^c\cap D_2\cap D_3\) instead have probabilities
\(p(1-p)^2\) and \(p^2(1-p)\).
Their difference is \(p(1-p)(1-2p)\), with the asserted
threshold at \(p=1/2\). No threshold statement for those
exclusive configurations changes the earlier inclusive
comparison or the source's explicitly inclusive meanings.
\end{proof}
\end{annexitem}

\begin{annexitem}{annex:p19.explanation-choice-model}{thm:p19-explanation-choice-model}{Response functions and alternative constructions}
\begin{proof}
For \(R>0\), interpret \(R^b=\exp(b\log R)>0\).
The defining denominator exceeds its numerator, so
\(0<q(R)<1\). Direct subtraction and division give
\[
 \frac{q(R)}{1-q(R)}=\exp(a)R^b,
 \qquad \log\frac{q(R)}{1-q(R)}=a+b\log R.
\]
For \(b=1\), \(q>1/2\) is equivalent to
\(\exp(a)R>1\). With \(R=5/8\) and multiplier four,
\[
 q=\frac{4(5/8)}{1+4(5/8)}=\frac57.
\]
The independently calculated event posteriors were
\(25/61,40/61\), neither \(5/7\).

For the threshold construction take a uniform variable
\(V\) on \((0,1)\), representing a population coordinate.
For \(b>0\), give person \(V=v\) the fixed threshold
\[
 T(v)=\frac{\log(v/(1-v))-a}{b}.
\]
That person deterministically chooses the simple answer
exactly when \(\log R\geq T(v)\). By multiplying by
positive \(b\), exponentiating, and solving for \(v\),
this is equivalent to
\[
 v\leq \frac{\exp(a)R^b}{1+\exp(a)R^b}.
\]
The uniform probability of this event equals \(q(R)\).
Every individual uses a deterministic rule, yet the
aggregate curve equals the proposed probability-matching
curve at every positive \(R\).
The construction asserts no distribution of actual
participants' thresholds; it proves non-uniqueness
of a psychological interpretation of the group curve.

For the causal-direction example take four equally weighted
fictional records. Give them pre-choice prevalence estimates
\(Z=.08,.08,.20,.20\), respectively. Set the choice indicator
\(A=\mathbf1_{\{Z>.15\}}\), where one means the simple
explanation, and define the later report \(Y=Z\).
Then \(\mathbb E[Y\mid A=1]=.20\) and
\(\mathbb E[Y\mid A=0]=.08\).
Both groups have probability one half, so these conditional
means have positive denominators.
The equation for \(Y\) does not use \(A\); even if the
choice value is forcibly replaced while keeping \(Z\)
fixed, the report remains \(Z\).
Thus a positive choice--report association and a report
elicited after choice can occur with no effect of choice
on the report. This is a specified countermodel, not
a claim that it generated the published data.
\end{proof}
\end{annexitem}

\begin{annexitem}{annex:p19.explanation-chi-tail}{thm:p19-explanation-chi-tail}{The printed chi-square reference calculation}
\begin{proof}
The gamma normalization and recurrence in
\cref{thm:p15-gamma-normalization}, together with the
Gaussian integral in \cref{thm:p15-normal-normalization},
give
\[
 \Gamma(1/2)=2\int_0^\infty e^{-u^2}\,du=\sqrt\pi,
 \qquad \Gamma(3/2)=\frac{\sqrt\pi}{2}.
\]
The substitution \(x=u^2\) is justified first on compact
positive intervals by differentiating the corresponding
integrals, then taking monotone limits at zero and infinity.
Therefore the declared shape-\(3/2\), rate-\(1/2\)
density is
\[
 f(t)=\frac{\sqrt t\,e^{-t/2}}{\sqrt{2\pi}}
 \qquad(t>0).
\]
For \(x>0\), the substitution \(t=2u^2\), justified by
the same finite-interval argument, yields
\[
 Q_3(x)=\frac4{\sqrt\pi}
             \int_{\sqrt{x/2}}^\infty u^2e^{-u^2}\,du.
\]
Integration by parts gives, for \(a>0\),
\[
 \int_a^\infty u^2e^{-u^2}\,du
 =\frac a2 e^{-a^2}
       +\frac12\int_a^\infty e^{-u^2}\,du.
\]
The exponential series gives \(e^{u^2}\geq u^4/2\) for
\(u>0\), so \(0\leq u e^{-u^2}\leq2/u^3\to0\).
This justifies the limiting endpoint.
Substitution of \(a=\sqrt{x/2}\) proves the stated
formula with the explicitly defined complementary
Gaussian integral.

Evaluating that formula at \(x=7.26\) gives
\(0.064056576\ldots\), which rounds to the handbook's
derived value .06406, not the source's printed \(p<.01\).
The strict comparison with .01 can also be proved
without relying on the precision of this numerical
evaluation. The positive density gives
\[
 Q_3(7.26)>Q_3(8)>
 \frac4{\sqrt\pi}e^{-4}>
 \frac2{81}>.01.
\]
Here the omitted complementary Gaussian integral is
positive, \(\pi<4\), and \(e<3\).
For completeness, \(\pi/4=\int_0^1(1+t^2)^{-1}dt<1\);
and the exponential series gives
\(e=2+\sum_{n\geq2}1/n!<2+\sum_{n\geq2}2^{-(n-1)}=3\).
Thus the discrepancy from \(p<.01\) is not a rounding
artifact in the displayed evaluation.

This calculation conditions on the printed value 7.26
and the printed three-degree reference law. It does not
identify the true contingency-table statistic or prove
that the authors actually used those inputs.
Neither an alternative observed statistic nor a participant
count table follows from this calculation.
\end{proof}
\end{annexitem}
